\section{Introduction}

A succinct non-interactive argument of knowledge (SNARK) allows a
prover to establish knowledge of a witness using a short proof that
can be verified much faster than the computation it
certifies~\cite{innovations/BitanskyCCT12}.
SNARKs enable efficient verification of computations performed by
untrusted parties; their zero-knowledge variants additionally
hide the witness (see, \eg,~\cite{tcc/BootleCG20,%
eurocrypt/BootleCL22,%
eurocrypt/BuenzFS20,%
eurocrypt/ChenBBZ23,%
eurocrypt/ChiesaHMMVW20,%
eurocrypt/ChiesaOS20,%
iacr/GabizonWC19,%
crypto/GolovnevLSTW23,%
eurocrypt/Groth16,%
crypto/Setty20,%
iacr/SettyL20}).
The prover's cost of producing these SNARK proofs, however, can substantially
exceed the cost of the computation itself.

The prover's use of a large working field is a major source
of this proving overhead.
Many efficient SNARKs combine a polynomial interactive oracle
proof (PIOP) with a polynomial commitment scheme
(PCS)~\cite{tcc/Ben-SassonCS16,eurocrypt/BunzFS20,%
eurocrypt/ChiesaHMMVW20,iacr/GabizonWC19,ccs/MallerBKM19}.
The circuit-satisfiability (Circ-SAT) relation is encoded as low-degree
polynomial identities.
In multilinear constructions, the PIOP uses random challenges
to reduce these identities to evaluation claims about the
multilinear extensions (MLEs) of the witness vectors, which
the PCS then authenticates.
To achieve the target soundness error, these protocols use
a challenge field sufficiently large relative to the polynomial
degrees, even when the witness entries fit in a much smaller field.
Arithmetic involving these challenges produces intermediate
values in the larger field, increasing both prover time
and memory usage.

Small fields and low-bit-width integers arise naturally in many modern
applications of SNARKs.
First, a growing body of zero-knowledge machine learning (zkML)
research generates SNARK proofs for neural-network (NN) computations
(\eg,~\cite{ccs/Balbas0VRS23,%
asiacrypt/CCYY25,%
nips/GhodsiGG17,%
iacr/GuoGQZ25,%
esorics/KeufferMC18,%
uss/Weng0XK021,%
ccs/LiuXZ21,%
iacr/FengQZDC21,%
corr/KangHSS22,%
ccs/GargGJMMPW23,%
tifs/WengWTYL23,%
tdsc/LeeKKO24,%
ccs/SunLZ24,%
ccs/AbbaszadehPK024,%
popets/WangH23,%
icml/WaiwitlikhitS0H24,%
sp/QuGYZ26,%
ccs/RiasiWBVH25,%
uss/GuoLHQTZ25,%
uss/0001SLLGCZ25}).
Modern production-grade NNs commonly use low-precision (\eg, 8-bit)
representations~\cite{corr/micikevicius22}.
Second, a recent line of work combines SNARKs with fully homomorphic
encryption (FHE) to achieve verifiable computation on encrypted
data (\eg, ~\cite{ccs/FioreGP14,%
pkc/FioreNP20,%
pkc/BoisCFK21,%
joc/GaneshNS23,%
cic/AtapoorBPS24,%
wahc/KnabenhansVMH24,%
crypto/GargGW24,%
asiacrypt/AranhaCGS25,%
crypto/CascudoCCFGS25,%
ccs/ThibaultW25,%
cic/GamaBKSV25,%
ccs/ZhangWLXDFL25,%
iacr/PeshawariaLFT25}).
FHE schemes often exploit small-field arithmetic, with plaintext
coordinates in fields such as
\(\mathbb{F}_{65537}\)~\cite{asiacrypt/AranhaCGS25}
and ciphertext coefficients represented via the Chinese remainder
theorem (CRT) over 30--60-bit prime fields.

For computations over small fields, we seek prover time and
memory usage that scale proportionally with the bit width of the base field
rather than that of the soundness field, a property we call
\emph{bitwidth awareness}.
In elliptic-curve-based PCSs such as KZG, Hyrax, and
Dory~\cite{asiacrypt/KateZG10,sp/WahbyTSTW18,tcc/Lee21},
short scalars reduce the number of group operations in the
multiscalar multiplications used to generate
commitments~\cite{eurocrypt/SettyTW24}.
Hash-based constructions exploit binary-field packing and
ring-switching to commit to small-field polynomials without
representing each input value as a full extension-field
element~\cite{eurocrypt/DiamondP25,eurocrypt/DiamondP26}.
Extending these savings to the PIOP layer is more difficult:
random challenges turn small-field inputs into extension-field
intermediate values.
Small-field-first algorithms, constraint packing, and related
sum-check optimizations substantially reduce this
overhead~\cite{iacr/BagadDT24,iacr/DaoT24,iacr/Gruen24,%
iacr/BagadDDT25,iacr/DaoDBDT26}.
After \(t\) optimized rounds, a conventional sum-check tail
operates on tables of length \(N/2^t\) and, for fixed product
degree \(d\geq2\), performs \(\Theta(N/2^t)\) extension-field
multiplications~\cite{crypto/Thaler13,iacr/BagadDT24}.
In a degree-\(\kappa\) extension represented in a polynomial basis,
schoolbook multiplication uses \(\Theta(\kappa^2)\) base-field
multiplications~\cite{iacr/BagadDT24}.
Increasing \(t\) reduces this tail at the expense of the work
in the optimized phase; \(t\) may depend on \(\kappa\).
The prover's total cost therefore depends on both phases.

Field agnosticism does not imply bitwidth awareness.
Field-agnostic constructions such as Brakedown and BaseFold
remove structural restrictions on the arithmetic field, but not
the field-size requirements imposed by
soundness~\cite{crypto/GolovnevLSTW23,crypto/ZeilbergerCF24}.
Their algebraic checks still rely on a sufficiently large
challenge field to achieve the target soundness error.
Supporting a small base field therefore does not by itself
ensure that challenge-dependent arithmetic and intermediate
storage retain the cost advantages of the low-bit-width witness~\cite{eurocrypt/DiamondP26,iacr/BagadDDT25}.
We therefore ask:
\begin{quote}
\emph{Can a SNARK achieve bitwidth-aware prover time and memory
usage while drawing its soundness challenges from a large field?}
\end{quote}

\smallskip\noindent\textbf{Our result.}
Let $\lambda$ be the security parameter, and consider an $N$-gate
arithmetic circuit over a small field of order $q$.
We choose a packing factor $\kappa\simeq\lambda/\log_2 q$,
subject to $\kappa\le q$; for fixed $\lambda$, $\kappa$ scales
inversely with the base-field bit width $\log_2 q$.
For a single instance, our prover uses
$O(\kappa N(\log_q N)^2)$ base-field operations and stores
$O(N\log_q N)$ base-field elements.
For a batch of $\kappa$ instances of the same circuit, the
amortized costs decrease to $O(\kappa N)$ base-field operations
and $O(N)$ base-field elements of memory per instance.
Under schoolbook arithmetic, multiplication in a degree-$\kappa$
extension costs $\Theta(\kappa^2)$ base-field operations, while
each extension-field element occupies $\kappa$ base-field
coordinates.
Thus, the amortized per-instance complexity corresponds to
$O(N/\kappa)$ extension-field multiplication equivalents and
the storage occupied by $O(N/\kappa)$ extension-field elements.

\smallskip\noindent{\bf Our core technique.}
Our \methodname framework is built on a novel technique called \emph{vector polynomial
packing}.
In the base-field instantiations, we partition a length-\(N=m\kappa\) vector over \(\Field\)
into \(\kappa\) blocks of length \(m\), and encode these blocks
as a length-\(m\) vector of polynomials using \(\kappa\)
linearly independent polynomials in \(\Field[X]\), each of
degree at most \(\kappa-1\).
This encoding expresses the original vector relation as a
low-degree polynomial identity involving length-\(m\) polynomial vectors.
Evaluating this identity at a uniformly random point
\(\ext{x}\in\ExtensionField\), where
\([\ExtensionField:\Field]=\kappa\), reduces it to a
length-\(m\) vector relation over \(\ExtensionField\).
For our inner- and Hadamard-product instantiations, an appropriate
choice of packing polynomials allows the prover to establish
the reduced claim using \(O(\kappa N)\) base-field operations.

We further show how to commit to a length-\(N\) vector over
\(\Field\) by packing its entries into a length-\(m\) vector
over \(\ExtensionField\) and committing to it using any
compatible PCS over \(\ExtensionField\).
Combining this packed PCS with our packed inner- and Hadamard-product protocol and
a bitwidth-aware permutation argument yields a bitwidth-aware
SNARK for the PLONK constraint system~\cite{iacr/GabizonWC19} of circuit satisfiability.


\smallskip\noindent{\bf Concrete performance.}
Our component benchmarks compare the packed reductions with a
small-field-first baseline. At \(N=2^{30}\) over an 8-bit base
field, the measured prover speedups are \(5.60\times\) for inner
products and \(4.31\times\) for Hadamard products, with a
\(4.48\times\) reduction in resident memory for both.
These measurements include packing and witness-dependent
preprocessing, and exclude commitments and PCS openings.
Section~\ref{sec:experiment} reports the full comparison across
8-, 16-, 32-, and 64-bit configurations and the measurement procedure.


\subsection{Related Work}
\label{sec:related-work}

\noindent{\bf PIOP--PCS framework for SNARKs.}
Mainstream SNARKs combine a polynomial interactive oracle
proof (PIOP) with a polynomial commitment scheme
(PCS)~\cite{eurocrypt/BunzFS20,eurocrypt/ChiesaHMMVW20}.
In this framework, circuit satisfiability is arithmetized as linear
constraints, such as rank-one constraint systems
(R1CS)~\cite{eurocrypt/GennaroGP013} or PLONK-style constraints~\cite{iacr/GabizonWC19,eurocrypt/ChenBBZ23},
which are then encoded as polynomial identities.
The PIOP uses random verifier challenges to reduce verification
of these identities to polynomial evaluation claims, and the
PCS authenticates the claimed evaluations against committed
polynomials~\cite{asiacrypt/KateZG10,tcc/PapamanthouST13}.
The Fiat--Shamir transform makes the resulting argument non-interactive in the
random-oracle model~\cite{crypto/FiatS86,tcc/Ben-SassonCS16}.

Early arithmetizations encoded constraints using univariate
polynomials~\cite{eurocrypt/GennaroGP013,cacm/ParnoHG016,iacr/GabizonWC19},
whereas subsequent works increasingly adopted multilinear extensions
(MLEs) over the Boolean hypercube~\cite{sp/WahbyTSTW18,crypto/XieZZPS19,crypto/Setty20,eurocrypt/ChenBBZ23}.
At the PIOP layer, polynomial identities arising from MLE-based
arithmetizations are typically reduced to sumcheck claims.
The sumcheck protocol was introduced by
Lund \etal~\cite{jacm/LundFKN92} and subsequently developed into a
fundamental building block for efficient proof
systems~\cite{jacm/GoldwasserKR15,crypto/Thaler13}.
A PCS then binds the prover to the oracle
polynomials and proves the evaluation claims produced by the PIOP.
Our work follows this PIOP--PCS framework and develops
bitwidth-aware techniques for the PIOP layer.

\medskip
\noindent{\bf Bitwidth-aware PCS constructions.}
Group-based PCSs~\cite{asiacrypt/KateZG10,tcc/PapamanthouST13,%
sp/BunzBBPWM18,sp/WahbyTSTW18,tcc/Lee21,asiacrypt/BunzMMTV21}
commit to polynomial coefficients or evaluation vectors using
multiscalar multiplications (MSMs).
MSM algorithms can exploit short integer scalars, making
commitment generation sensitive to input bit
width~\cite{eurocrypt/SettyTW24}.
Opening costs depend on the construction: for a size-$N$ polynomial, Dory requires
\(O(\sqrt{N})\) group operations and \(O(N)\) scalar-field
operations~\cite{tcc/Lee21,eurocrypt/SettyTW24}, where
multiplications between short witness entries and full-width
weights can benefit from bitwidth-aware arithmetic.

Hash-based PCSs~\cite{icalp/Ben-SassonBHR18,crypto/ArnonCFY24,%
eurocrypt/ArnonCFY25,iacr/KohritaT23,ccs/AmesHIYM17,%
focs/Ben-SassonCIKS20,crypto/GolovnevLSTW23,crypto/XieZS22,%
crypto/BlockFK24,iacr/LiLXYY26,crypto/ZeilbergerCF24,%
usenix/GuoLHQTZ25,eurocrypt/BrehmCFRRZ25,%
asiacrypt/NariST25,iacr/BunzFRW25}
commit to linear encodings of polynomial coefficient or
evaluation vectors using Merkle trees.
Their opening protocols combine algebraic consistency checks
with authenticated queries to the encoded vectors, using
code distance to detect inconsistencies.
When encoding is performed over a large field, short inputs
can produce full-width codeword symbols, so commitment
generation does not automatically benefit from low input
bit width.
Recent constructions avoid this large-field embedding
overhead for small-field polynomials.
Binius~\cite{eurocrypt/DiamondP25} adapts Brakedown to
multilinear commitments over binary towers without
coefficient-wise embedding overhead.
FRI-Binius~\cite{iacr/BenjaminDP24} combines the binary tower of Binius with binary-field FRI and BaseFold,
achieving polylogarithmic proof size and verifier complexity.
Ring-switching generalizes this approach by reducing small-field
multilinear openings to openings of packed polynomials over
an extension field.
By packing \(\kappa\) base-field coordinates into each element
of a degree-\(\kappa\) extension, it transforms an extension-field
multilinear PCS into a small-field PCS whose commitment cost
matches that of the underlying scheme on inputs of the same
bit size.

% Complete degree-2 Hadamard audit: a_i b_i - c_i = 0.
% N is the length of EACH of the three input vectors (total input Theta(N)).
% The public equality factor is included but is not a third small-field input.
%
% SFF: 2025/1117, Theorem 6.1/Section 6.2; 2026/587, Figure 2,
% Section 6 and Appendix C.5. With M=N/2^t, the quadratic grid has 3^t
% entries (projective nodes 0,1,infinity). M*3^t mixed products cost
% kappa*N*(3/2)^t; suffix weights and tail cost kappa^2*M;
% grid contractions cost kappa^2*3^t. Streamed binding is absorbed.
% Gruen: 2024/108, Sections 2.3.1 and 5.1--6.1; 2026/587 Section 7.
% The full bound ADDS FFT interpolation N*t, equality weighting and
% binding kappa*N, suffix weights/tail kappa^2*N/2^t, and first-round
% Lagrange/quotient processing kappa^2*2^t. It is not the abstract's
% constraint-evaluation-only bound. This count uses FFT/coset domains;
% the protocol permits general interpolation sets (Section 5.1, footnote 18).
% Flock: 2026/1329, Sections 4.2--4.4 (July 31, 2026 PDF).
% Byte chunks/subfield degree 8; 3<=t<=7, implementation t=6, kappa=128.
% Re-cost lookup LDEs, weighting, binding, output check, and conventional
% tail. Full time adds kappa^2*2^t; RAM adds kappa*2^t (fixed tables
% included). N>=4^t absorbs these terms. This is a derived bound.
% Dao--Thaler: 2024/1038, Sections 7--8, Algorithm 2, composed here with
% CONVENTIONAL sumcheck. This is not a claim about the best possible tail.
% Liu--Zhang: 2025/594, Section 4.2, Theorem 8/Protocol 5 (NOT Table 1).
% Set s=2, h=floor(log_2(kappa)/2), u_i=4^i for i<h.
% Full cost: sum_i N/2^i*(u_i^2+kappa*u_i) + kappa^2*N/2^h
% = O(kappa^(3/2)*N). Generate deterministic prefix weights using the
% O(kappa) special-generator multiplication of 2024/1038 Section 8.
% Store only random-suffix weights and the current folded witness tables:
% O(N + max_i N*2^i + kappa*N/2^h) = O(N*sqrt(kappa)), for N>=kappa.
% RMFE: 2021/694, Theorem 4/Figure 3/Appendix A.2, sparse Boolean R1CS:
% n=3N, m=N, nnz=Theta(N), constant-rate packing and security O(kappa).
% Packed sumcheck: 2025/719 (2025-10-03 version), Lemma 5 explicitly proves
% the Hadamard relation via Lemma 2 (d=2); Figure 3 includes the output c.
% Set r=k, N=k*N_0. Equality-weighted degree is 3, hence q>3(2r-1).
% Stream per-column forward differences; keep O(N+r^2) base coordinates.
% Our row adds the output-vector linear evaluation to Proposition
% proposition:hadamard-complexity, costing O(kappa*N) and O(N+kappa^2).
% Parameter-only additive terms and their absorption conditions are given
% in Appendix appx:hadamard-comparison.
\begin{table*}[!t]
\centering
\setlength{\tabcolsep}{1.5pt}
\renewcommand{\arraystretch}{1.15}
\begin{threeparttable}
\caption{Prover costs for length-\(N\) Hadamard-product proofs}
\label{tab:hadamard-comparison}
\begin{tabular*}{\textwidth}{@{\extracolsep{\fill}}
    >{\raggedright\arraybackslash}p{0.25\textwidth}
    >{\raggedright\arraybackslash}p{0.26\textwidth}
    >{\raggedright\arraybackslash}p{0.20\textwidth}
    >{\raggedright\arraybackslash}p{0.25\textwidth}@{}}
\toprule
Method & Online prover time & RAM
& Applicability \\
\midrule
Classical sumcheck
& \(O(\kappa^2N)\)
& \(O(\kappa N)\)
& General extensions \\
\addlinespace[3pt]
SSF~\cite{iacr/BagadDT24,iacr/BagadDDT25,iacr/DaoDBDT26}
& \(O(\kappa N(3/2)^t +\kappa^2N/2^t)\)
& \(O(N+\kappa N/2^t)\)
& General extensions \\
\addlinespace[3pt]
Dao--Thaler~\cite{iacr/DaoT24}
& \(O(\kappa^2N)\)
& \(O(\kappa N)\)
& Binary tower fields \\
\addlinespace[3pt]
Tower zerocheck~\cite{iacr/LiuZ25}
& \(O(\kappa^{3/2}N)\)
& \(O(N\sqrt{\kappa})\)
& Binary tower fields \\
\addlinespace[3pt]
RMFE--Aurora~\cite{iacr/CascudoG21}
& \(O(\kappa^2N+\kappa N\log N)\)
& \(O(\kappa N)\)
& Binary RMFE \\
\addlinespace[3pt]
Packed sumcheck~\cite{iacr/WeiWXZDWZ25}
& \(O(rN)\)
& \(O(N+r^2)\)
& \(|\Field|>3(2r-1)\) \\
\addlinespace[3pt]
Gruen~\cite{iacr/Gruen24}
& \(O(Nt+\kappa N +\kappa^2N/2^t)\)
& \(O(N+\kappa N/2^t)\)
& $|\Field| \ge 2^{t+1}-1$ \\
\addlinespace[3pt]
Flock~\cite{iacr/BunzRW26}
& \(O(N2^t+\kappa N+\kappa^2N/2^t)\)
& \(O(N+\kappa N/2^t)\)
& \(\Field=\mathbb F_2\) \\
\midrule
\methodname
& \(O(\kappa N)\)
& \(O(N)\)
& General extensions \\
\bottomrule
\end{tabular*}
\begin{tablenotes}[para,flushleft]
\footnotesize
\item Here \(\kappa=[\ExtensionField:\Field]\) and
\(t_{\mathbb K}=[\mathbb K:\Field]\) is chosen as in
Section~\ref{sec:hadamard}.
Arithmetic and RAM are measured in base-field operations
and elements, respectively.
\(t\) counts optimized
or skipped rounds, and \(r\) is the packed-sumcheck parameter.
Bounds use schoolbook arithmetic and coordinate-charged table accesses;
RAM includes inputs and read-only tables. Public preprocessing is separate.
Parameter-only additive terms are absorbed for sufficiently large \(N\):
\(N\geq6^t\) for SFF and \(N\geq4^t\) for univariate skip and Flock.
Both \methodname{} rows assume \(N\geq\kappa^2\).
Flock uses \(\mathbb F_{256}\subseteq\ExtensionField\) and
\(3\leq t\leq7\) (\(t=6\) in its implementation).
Complete bounds, implementation assumptions, and soundness-dependent
parameter choices are given in Appendix~\ref{appx:hadamard-comparison}.
\end{tablenotes}
\end{threeparttable}
\end{table*}

\medskip
\noindent{\bf Low-bit-width PIOP optimizations.}
In sumcheck-based PIOPs, random challenges produce
extension-field intermediate values even when the witnesses
lie in a small base field.
Existing approaches reduce this overhead in several ways.

Small-field-first methods precompute products over the base
field and postpone extension-field arithmetic, but switching
to conventional sumcheck after \(t\) optimized rounds still
leaves \(\Theta(N/2^t)\) extension-field multiplications for
fixed product degree \(d\geq 2\), while increasing \(t\)
requires more base-field
preprocessing~\cite{iacr/BagadDT24,iacr/BagadDDT25,iacr/DaoDBDT26}.
Gruen's univariate-skipping technique replaces several initial
rounds with one higher-degree check, reducing prover work
at the cost of higher-degree verifier
interpolation~\cite{iacr/Gruen24}.
Flock applies univariate skipping to batched Boolean computations
and computes binary-input low-degree extensions using chunked
lookup tables~\cite{iacr/BunzRW26}. It compresses these tables
using translation invariance, improving cache locality without
changing the number of vector XORs. Its zerocheck also uses
fixed equality weights and aggregates the linear output term
before interpolation.

Binary-tower constraint packing combines base-field constraints
using linearly independent extension-field
weights~\cite{eurocrypt/DiamondP25,iacr/DaoT24}.
Dao and Thaler~\cite{iacr/DaoT24} use inexpensive multiplication
by tower generators to reduce equality-table construction to
approximately \(N/\kappa\) general extension-field multiplications,
but their direct variable elimination requires individual
degree one and therefore does not yield the same bound for
proving general quadratic relations.
Liu and Zhang~\cite{iacr/LiuZ25} specialize sumcheck to Boolean
inputs over binary towers, achieving \(O(\kappa^{3/2}N)\)
prover time under quadratic multiplication costs, but this
bound relies on the binary tower structure over
\(\mathbb{F}_2\).


Reverse multiplication-friendly embeddings allow coordinatewise
Boolean products to be recovered from packed extension-field
products, but the resulting Boolean R1CS protocols also require
specialized checks of linear relations modulo a
subspace~\cite{iacr/CascudoG21}.

Packed sumcheck instead keeps computation in a small prime
field by combining parallel repetition with folding, but
the larger soundness error of individual small-field checks
requires additional work for soundness amplification and
for combining the repeated
claims~\cite{iacr/WeiWXZDWZ25}.


\medskip
\noindent{\bf Relation to our work.}
Our work combines packing at the relation and commitment layers.
For inner and Hadamard products, vector polynomial packing
encodes operand blocks as row polynomials and uses auxiliary
polynomials to account for cross-column products.
In the base-field instantiations, the prover computes row
products over \(\Field\), then evaluates at a random point in
\(\ExtensionField\) to reduce
the length-\(N\) claim to a length-\(N/\kappa\) inner-product
claim.
Our packed opening protocol authenticates these
challenge-dependent operands against a fixed-basis witness
commitment.
Together with a packed permutation argument, these components
support PLONK constraints while maintaining compact witness
representations.
Sections~\ref{sec:overview} and~\ref{sec:pcs} describe the
reductions and their commitment interface.
For \(|\Field|<\kappa\) and power-of-two \(\kappa\), our
univariate Hadamard construction
uses an intermediate field and interpolates only after weighted
aggregation. Combining blockwise linear-map tables, as used for
binary LDEs in Flock~\cite{iacr/BunzRW26}, with intermediate-field
multiplication tables yields \(O(N\kappa+\kappa^3)\) online
base-field operations. The construction and its table costs
are given in Section~\ref{sec:hadamard}.


\subsection{Our Contributions}

\noindent\textbf{Vector polynomial packing.}
    We introduce a technique that reduces length-\(N=m\kappa\)
    inner- and Hadamard-product relations over \(\Field\) to
    length-\(m\) inner-product relations over a degree-\(\kappa\)
    extension \(\ExtensionField\).
    In the base-field instantiations, the prover multiplies row polynomials over \(\Field\) before
    evaluating them at an extension-field challenge.
    This keeps the polynomial-product computation in the small field
    and confines the ensuing extension-field proof to vectors shorter
    by a factor of \(\kappa\).

\medskip \noindent \textbf{End-to-end bitwidth-aware construction.}
    We compose our packed reductions and permutation argument with
    a compatible bitwidth-aware PCS through ring-switching-based
    openings to obtain arguments for PLONK-style circuit satisfiability.
    We give constructions for both a single instance and a batch of
    instances of the same circuit.
    Packing across instances preserves the shared linear constraints,
    yielding lower amortized prover costs.
    Section~\ref{sec:together} gives soundness and efficiency bounds
    for both compositions, including their dependence on the
    underlying PCS.

\medskip\noindent \textbf{Independently useful Hadamard-product and permutation proofs.}
    Our bitwidth-aware protocols for coordinatewise multiplication
    and multiset equality apply beyond the circuit arithmetization
    considered here.
    For \(q=|\Field|<\kappa\), a tabulated intermediate-field
    Hadamard reduction uses \(O(N\kappa)\) online base-field
    operations and \(O(N)\) working base-field coordinates,
    assuming \(\kappa\) is a power of two and \(N\geq\kappa^2\).
    It retains the original witness commitments; its reusable
    tables and their preprocessing costs are accounted for
    separately in Section~\ref{sec:hadamard}.
    The permutation proof uses packed Hadamard checks to verify
    grand-product identities while retaining compact witness storage.
    It also provides a reusable component for permutation-based
    lookup arguments.
